\subsection{Threshold bounds}
We use exactly the parameters and menus of Section~\ref{sec:refined-mechanism}.
Put $K=A-q=721/1500$, $\rho_*=b-A=165649/500000$ and
$B_0(k)=C_0(k)-k$. On constrained $Q$, $c\le k\le K<d$ and
\[
 d<B_0(K)\le B_0(k)\le B_0(c)<A,\qquad
 B_0(K)-h(\rho_*)=\frac{37}{5000000}>0.
\]
The function $B_0$ decreases on this interval. Outside the fee support,
$t-k\ge q$ gives $z\ge k+h(\rho)$. On its support, $Q$ implies
$k\le\min(K,b-\rho-q)$. For $\rho\le\rho_*$, the displayed slack
bounds $B_0(k)-h(\rho)$ below. For $\rho\ge\rho_*$, its sufficient
lower bound $B_0(b-\rho-q)-h(\rho)$ has derivative at least
$\lambda-3K/2=79/1000>0$, so the same bound applies. Consequently
\[
 C=\max(C_0(k),z)
\]
on every constrained $Q$ report. If $B$ is capped at $A$, then $C=z>b_0$.
This argument includes the fee face and does not assume $C_0(k)\le1$.

A second bound needed below is
\[
 C_0(k)-k-d\ge g(k).
\]
Off the support it follows from $B_0(k)>d$. On $[c,\bar u]$, the
difference has positive derivative $3k/2-1+\lambda$ and minimum
$212401/3000000$. Since $t\le h(k)$, this also proves $C_0(k)\ge t+c$.
The inverse implications $x<k(t)\Rightarrow h(x)<t$ and
$x\ge k(t)\Rightarrow h(x)\ge t$ hold through the whole plateau;
an identity $h(k(t))=t$ is not used.

\subsection{Base/base compatibility}
Before the fee, the base menus are the pivot-normalized menus of a common
nine-outcome affine maximizer. Its costs are zero for empty, $a$ for either
singleton allocated to one bidder, $b$ for either bidder's bundle, and $s$
for either split allocation. Selecting a maximizer of total reported value
minus cost supplies a jointly feasible pair of menu maximizers. Deleting
any zero-utility allocation preserves such a pair with the stated empty rule.

The prices satisfy the strict bundle-discount inequality
\[
 P_{\rm scarce}+P_{\rm safe}-P_B\ge q>0.
\]
If both opponent coordinates are at most $d$, the left side is at least
$a-c>q$. If exactly one exceeds $d$, use $H\ge t-a$ to obtain $q$;
if both exceed $d$, use $H\ge z-b$ to obtain $2q$.

Adding the same nonnegative fee to safe and bundle preserves a maximizing
subset of every uncharged maximizer. Indeed, if safe maximizes, comparison
with bundle gives $x\le P_B-P_{\rm safe}<P_{\rm scarce}$ for the own
scarce value $x$. The scarce singleton has negative utility, and the
bundle-minus-safe comparison is unchanged by the fee, leaving safe or empty
as a maximizer. A scarce or empty maximizer remains a maximizer because
its price is unchanged. A bundle contains every possible new allocation.
Applying this argument to the jointly feasible uncharged pair proves
base/base compatibility, including positive ties.

\subsection{Q/base and Q/Q compatibility}
Fix a $Q$ report $w=(t,\rho)$, with physical item 1 high, and write the
other report as $v=(x,y)$. A base singleton costs at least $d$ and its bundle
at least $b$. Since $t+\rho\le b$, an own type $w$ can buy only its high
item positively in a base menu; if $t\le d$ it must be empty.

For $y\le d$, its item-1 base price is at least $\max(x,x+y-c)$.
An opposing $Q$ scarce purchase with positive utility requires $x>A\ge t$.
A selected opposing bundle has $x+y\ge C\ge C_0(k)\ge t+c$.
Whenever the $Q$ bidder consumes item 1, the base high price is therefore
at least $t$, so the base bidder has no positive utility and selects empty.

For $y>d$, the base price is at least $x+q$. A positive high purchase by
$w$ would require $x<t-q\le K<d<y$. Thus $x$ is the low coordinate of
the base report and the fee on that physical item gives a price at least
$h(x)$. The opposing $Q$ scarce singleton requires $x\ge A>k$;
a bundle requires $x\ge C-B\ge k$ by comparison with safe. In either
case $h(x)\ge t$, contradicting the positive base purchase. The inequalities
include the inverse knots, the fee face and capped bundle rows.

For $Q/Q$, each own coordinate is at most $A$. Neither a scarce singleton
at $A$ nor a capped safe singleton at $A$ has positive utility. An uncapped
singleton costs more than $d$, while a constrained $Q$ type's low coordinate
is below $d$. Positive singletons can therefore only be distinct own high
items under opposite orientations. Two positive bundles are impossible,
because each bundle price is at least the opposing report's total value.

For a possible bundle/safe conflict under opposite orientations, a positive
safe buyer with high value $y$ and opposing low value $\rho$ requires
$y>B\ge h(\rho)$. A conflicting bundle in the reverse menu requires
$\rho\ge k(y)$; the cap only increases this upgrade threshold. Hence
$h(\rho)\ge y$, a contradiction. The strictness comes from the safe
buyer's positive utility, so a tied bundle/safe comparison on the other
side is covered. Same orientations have no positive singleton. Free $Q$
types cannot buy either singleton positively and cannot buy the opposite
constrained safe option positively. Empty at zero completes the tie cases.

\subsection{Compatibility of lottery menus}
For an $E$ menu put $L=\delta+\alpha/2$, $Y=c+L$ and $j=1-t/2$.
The factored formula for $\delta$ gives, throughout $A<t<T$,
\[
 \beta L=\alpha/2,\quad t+c+\delta-Y=j,\quad
 c<Y<d<d+\delta<A,\quad d<j<A<t,
\]
with $0<\beta<1$ and bundle upgrade threshold $t-q>c$.

Consider $E$ against an uncharged base menu first. An own $E$ type
$w=(t,\rho)$, $\rho<c$, cannot buy its low base singleton; its base
bundle is dominated by high because $P_B-P_1\ge c>\rho$. Only item 1
can conflict. If the other report $(x,y)$ has $y\le d$, the base high
price is at least $\max(x,x+y-c)$. An $E$ scarce purchase requires
$x\ge t$, a bundle requires $x+y\ge t+c+\delta$, and a lottery requires
$x+\beta(y-c)\ge t$. For a lottery, $y\ge c$ implies $x+y-c\ge t$,
while $y<c$ implies $x\ge t$. Thus every allocation consuming item 1
makes the opposing base bidder empty. If $y>d$, the base high price is
at least $x+q$, the lottery cannot win because $y>Y$, and a bundle
requires $x\ge t-q$. Again the base bidder is empty. The maximizing-subset
argument permits the actual base fee without changing the $E$ allocation.

For $E/Q$, an own $Q$ report has coordinates at most $A$ and sum at most
$b=a+c$. Its $E$ scarce singleton has negative utility. Its lottery utility
is at most $A-t+\beta(a-A)$: for safe value $y\ge c$, write the lottery
value as $(1-\beta)x+\beta(x+y-c)$; for $y<c$, scarce dominates lottery.
Now
\[
 (t-A)/\beta=t-A+2\delta/3
\]
is strictly increasing, with derivative at least $d+c>0$ and infimum
$q^2/2=a-A$ at $A$. The lottery utility is strictly negative on the open
$E$ interval. Bundle utility is at most $a-t-\delta<0$, because
$t-A+\delta$ increases and $\delta(A)=3q^2/4>a-A$.

Only safe can be purchased positively. Its value $y>d$ is the own $Q$
high coordinate. The reverse $Q$ menu has scarce-upgrade threshold at least
$k(y)\ge c$, so its own $E$ low value $\rho<c$ cannot obtain the
conflicting item as scarce or bundle. This covers the inverse plateau using
the weak inequality $k(y)\ge c$.

For $E/E$ with the same physical orientation, an own low value below $c$
makes bundle and lottery dominated by scarce. Only the larger high value
can buy that item; equal high values give zero utility. Under opposite
orientations, each scarce value is below $c$, the bundle upgrade $t-q$ and
the lottery cell's minimum scarce value $j>d$. Only the disjoint own high
items can be allocated. These cases exhaust all menu pairs. Closed $Q$,
open $E$ and the remaining base rule cover all report boundaries.

\subsection{Exact integration of the seed}
For a fixed opponent report $(t,\rho)$ on
$\mathcal T=\{0\le\rho\le t\le1\}$, each menu option has an affine
utility in the own type. Its demand region is the intersection of its
utility-comparison half-planes with the unit square. Therefore its payment
contribution is its price times that polygon's area; lottery probabilities
enter both the half-planes and the item-count integrand. With $r_c$ as
defined in \eqref{eq:v5-face-reductions},
\[
 R_*=4\int_{\mathcal T}r_c(t,\rho)\,dt\,d\rho.
\]
The factor four is two bidders and two opponent orientations. The integral
includes the entire inverse plateau, the strip $A<t<a$ in $E$, and the
fee strip $c\le\rho\le b-A$, $A<t<b-\rho$.

Exact polynomial/simplex integration and a separate decomposition using
Green boundary integrals agree on all three algebraic coefficients in
\eqref{eq:refined-revenue}. An implementation importing neither calculator
reconstructs the full integral by vertical slicing and obtains the same
coefficients. These are complete evaluations of the stated seed menus;
no revenue from a different parameter choice is added. Positive-coordinate
tie interfaces have zero own-type area, while the pointwise argument above
defines the mechanism on those interfaces. The package README identifies
the structural proof and each independent integral implementation.

\subsection{Certified reserve choice and family boundary}\label{app:reserve-parameter}
For the fixed seed, the face formula is a polynomial $R(t)$ of degree
at most five. Write it as $\sum_m m w_m(t)$, where the twelve fixed moments
$m$ have rational enclosures and $w_m$ are the explicit polynomial weights
in \eqref{eq:v5-revenue}. To bound $R(t_1)-R(t_2)$, form
$w_m(t_1)-w_m(t_2)$ first and then choose the appropriate endpoint of each
moment interval according to that coefficient's sign. This retains the
common-moment dependence instead of subtracting two independent revenue intervals.
It gives
\[
 R(83/10000)-R(1/100)>0.0000021293,
\]
while
\[
 -0.002506136<R'(1/100)<-0.002504465<0.
\]

The family optimum can also be localized exactly. If a polynomial on
$[a,b]$ is expanded as $p(a+(b-a)z)=\sum_{k=0}^n c_kz^k$, its Bernstein
coefficients are
\[
 b_i=\sum_{k=0}^i c_k\frac{\binom{i}{k}}{\binom{n}{k}},
 \qquad 0\le i\le n.
\]
Each coefficient here is a linear form in the same fixed moments. Exact
interval bounds for every coefficient prove
\[
\begin{aligned}
 R'(t)&>0 &&(0\le t\le0.008297),\\
 R'(t)&<0 &&(0.008299\le t\le1),\\
 R''(t)&<0 &&(0.008297\le t\le0.008299).
\end{aligned}
\]
The polynomial extends continuously to $R(1)=0$, although the mechanism
is defined only for $t<1$. The derivative signs and strict concavity on
the middle interval imply a unique global family maximizer
\[
 0.008297<\tau_*<0.008299.
\]
Bernstein bounds on that middle interval also give
$\max_{0\le t<1}R(t)<0.87651436$. This is a fixed-seed family ceiling,
not an unrestricted auction upper. The selected simple rational
$\tau=83/10000$ is near the family maximizer and is not asserted to be
family-optimal or globally optimal.

The complete outward-rounded endpoints used by the release are
\[
\begin{aligned}
 R_\tau&\in[0.876514341027067549900397423113,\\[-2pt]
       &\hspace{2.5em}0.876514355424052828335888359874],\\
 U&\in[0.882351585488796576205586892263,\\[-2pt]
       &\hspace{2.5em}0.882351585488796576205586892264],\\
 U-R_\tau&\in[0.005837230064743747869698532390,\\[-2pt]
       &\hspace{2.5em}0.005837244461729026305189469151].
\end{aligned}
\]
The endpoint strings record rational bounds; their printed length is not
the precision of the unknown revenue. The parameter checker uses only
standard-library rational arithmetic and imports no face calculator.
Its input hashes, all Bernstein coefficient intervals, exact gain and
updated ownership-box margins are in the reserve-parameter manifest.
Finite implementation checks cover reserve faces, inverse knots and menu
boundaries separately from the all-real preservation theorem.
